Finden Sie eine Lösung der Differentialgleichung
\[
zy''(z) + by'(z) - y(z) = 0
\]
mit der Anfangsbedingung $y(0) = 1$ als Potenzreihe.

\begin{loesung}
Wir verwenden einen Ansatz als Potenzreihe für die Funktion $y(z)$
und berechnen die Ableitungen
\[
\renewcommand{\arraycolsep}{1pt}
\begin{array}{rcrcrcrcrcrcrcrcl}
   y(z) &=&  a_0 &+&        a_1z &+&         a_2z^2 &+&         a_3z^3 &+&         a_4z^4 &+&        a_5z^5 &+&           a_kz^k &+& \dots\\
 by'(z) &=& ba_1 &+&      2ba_2z &+&       3ba_3z^2 &+&       4ba_4z^3 &+&       5ba_5z^4 &+&      6ba_6z^5 &+& (k+1)ba_{k+1}z^k &+& \dots\\
zy''(z) &=&      & & 2\cdot1a_2z &+& 3\cdot 2a_3z^2 &+& 4\cdot 3a_2z^3 &+& 5\cdot 4a_5z^4 &+& 6\cdot5a_6z^5 &+& (k+1)ka_{k+1}z^k &+& \dots
\end{array}
\]
Die Differentialgleichung kann man auch in der Form
\[
zy''(z) + by'(z) = y(z)
\]
schreiben.
Setzt man die Potenzreihen in diese Form der Gleichung ein, erhält man die
Rekursionsgleichung
\[
ba_1=a_0
\qquad\text{und}\qquad
(k+1)(b+k)a_{k+1} = a_k
\]
für $k>0$.
Durch Auflösen nach $a_{k+1}$ erhalten wir
\[
a_{k+1}
=
\frac{1}{(b+k)(k+1)} a_k,
\]
welche für $k=0$ auch mit der unabhängig gewonnenen Bedingung
$a_1=a_0/b$ übereinstimmt.
Daraus kann man ableiten
\[
a_k
=
\frac{1}{b(b+1)\cdots(b+k-1)\cdot k!} a_0.
\]
Der erste Teil des Nenners ist das sogenannte Pochhammer-Symbol
\[
(b)_k
=
b(b+1)(b+2)\cdots (b+k-1),
\]
damit kann man den Koeffizienten $a_k$ schreiben als
\[
a_k = \frac{1}{(b)_k\cdot k!}a_0.
\]
Die Lösungsfunktion ist daher
\[
y(z)
=
a_0
\sum_{k=0}^\infty
\frac{1}{(b)_k\cdot k!} z^k.
\]
Die Reihe auf der rechten Seite ist bekannt als die hypergeometrische
Funktion
\[
\prescript{}{0}F_1\biggl(\begin{matrix}\text{---}\\b\end{matrix};z\biggr)
=
\prescript{}{0}F_1\Bigl(\text{---};b;z\Bigr)
=
\sum_{k=0}^\infty
\frac{1}{(b)_k\cdot k!} z^k.
\]
Die Lösung der Differentialgleichung ist daher die Funktion
$ y(z) = a_0 \prescript{}{0}F_1(\text{---};b;z)$.
\end{loesung}
